\section*{अध्याय \ref{ch:b2:microbial-metabolism} --- सूक्ष्मजैविक उपापचय और जैवफ़िल्में}

\begin{solution}{exo:b2:microbial-metabolism:1}
\begingroup\sloppy
सायनोजीवाणु: प्रकाश-शिला-स्वपोषी (प्रकाश, जल, $\mathrm{CO_2}$)।
\emph{Nitrosomonas}: रस-शिला-स्वपोषी (अमोनियम का ऑक्सीकरण,
$\mathrm{CO_2}$)। ग्लूकोज़ पर \emph{E.~coli}: रस-कार्बनिक-परपोषी।
प्रकाश में सक्सिनेट पर बैंगनी अ-गंधक जीवाणु:
प्रकाश-कार्बनिक-परपोषी। अँधेरे में सल्फाइड पर \emph{Thiobacillus}:
रस-शिला-स्वपोषी।
\par\endgroup
\end{solution}

\begin{solution}{exo:b2:microbial-metabolism:2}
ऑक्सीजन (सतह के मिलिमीटर), नाइट्रेट (ठीक नीचे, जहाँ ऑक्सीजन चुक गई
हो), मैंगनीज़ और लोहे के ऑक्साइड (भूरे से धूसर का संक्रमण), सल्फेट
(काली सल्फाइडी परत, समुद्री अवसादों में मोटी), \omterm{def:b2:microbial-metabolism:anaerobic}{मेथेनजनन} के लिए
$\mathrm{CO_2}$ (सबसे गहरे, या सल्फेट-दरिद्र मीठे जल में सतह के ठीक
नीचे)।
\end{solution}

\begin{solution}{exo:b2:microbial-metabolism:3}
$5\,\mathrm{C_6H_{12}O_6} + 24\,\mathrm{NO_3^-} + 24\,\mathrm{H^+}
\to 30\,\mathrm{CO_2} + 12\,\mathrm{N_2} + 42\,\mathrm{H_2O}$: प्रति
ग्लूकोज़ 4.8 नाइट्रेट आयन (हर नाइट्रेट 5 इलेक्ट्रॉन लेता है, हर
ग्लूकोज़ 24 देता है)।
\end{solution}

\begin{solution}{exo:b2:microbial-metabolism:4}
सतह से जुड़ा ऐसा सूक्ष्मजैविक समुदाय जो बहुशर्कराओं, प्रोटीनों और DNA
की स्वयं बनाई आधात्री में धँसा हो। चरण: संलग्नन, सूक्ष्मनिवह और
आधात्री-स्राव, मीनारों तथा नलिकाओं सहित परिपक्वन, प्रकीर्णन। दंत
पट्टिका और दंतक्षय; कैथेटरों तथा प्रत्यारोपों के (और सिस्टिक फ़ाइब्रोसिस
में फेफड़ों के) संक्रमण; पाइपों और जहाज़ों के पेंदों का जमाव तथा
संक्षारण --- या, उपयोगी रूप में, जल-उपचार के रिसाव-निस्यंदक।
\end{solution}

\begin{solution}{exo:b2:microbial-metabolism:5}
आठ इलेक्ट्रॉन \qty{-0.42}{V} से \qty{-0.24}{V} तक गिरते हैं:
$\Delta G^{\circ\prime} = -8\times 96.5\times 0.18 =
\qty{-139}{kJ/mol}$ प्रति मेथेन (सारणीबद्ध \qty{-131}{kJ/mol}), यानी
प्रति $\mathrm{H_2}$ \qty{35}{kJ}: अधिक से अधिक 0.7 ATP प्रति
हाइड्रोजन, व्यवहार में लगभग आधी। प्रति क्रियाधार-अणु ऊर्जा इतनी कम है
कि मेथेनजनक धीरे बढ़ते हैं (घंटों से दिनों के \omterm{def:b2:unicellular-diversity:division}{पीढ़ी काल}) और
अपना लगभग सारा क्रियाधार कोशिकाओं के बजाय गैस में बदल देते हैं।
\end{solution}

\begin{solution}{exo:b2:microbial-metabolism:6}
$\mu = 1.2\,S/(5 + S)$: \qty{0.20}{h^{-1}} ($T = \ln 2/\mu =
\qty{3.5}{h}$), \qty{0.60}{h^{-1}} (\qty{1.2}{h}), \qty{1.09}{h^{-1}}
(\qty{0.64}{h})। नब्बे प्रतिशत: $S = 9K_s = \qty{45}{mg/L}$।
\end{solution}

\begin{solution}{exo:b2:microbial-metabolism:7}
$S^* = 0.05\times 0.4/(0.8 - 0.4) = \qty{0.05}{g/L}$; $X^* =
0.4\times(5 - 0.05) = \qty{1.98}{g/L}$; निर्गत $DX^* = 0.4\times
1.98 = \qty{0.79}{g/L/h}$; बहाव $D_c = 0.8\times 5/5.05 =
\qty{0.79}{h^{-1}}$।
\end{solution}

\begin{solution}{exo:b2:microbial-metabolism:8}
पकड़ी गई ऊर्जा: $0.25\times 275 = \qty{69}{kJ}$, यानी प्रति अमोनियम
1.4 ATP; $12/1.4 = 8.7$, यानी प्रति कार्बन लगभग नौ अमोनियम आयन --- और
नाइट्राइट से इलेक्ट्रॉन चढ़ाकर NADH बनाने की लागत गिनें तो कई गुना
अधिक। किसी परपोषी को उसका कार्बन पहले से अपचयित और सुसंगठित मिलता है,
और प्रति ग्लूकोज़ तीस ATP; किसी नाइट्रीकारी को अपनी ऊर्जा और अपनी
अपचायक शक्ति दोनों किसी घटिया दाता से ख़रीदनी पड़ती हैं और हर कार्बन
$\mathrm{CO_2}$ से गढ़ना पड़ता है, इसलिए वह जितना क्रियाधार निपटाता है
उसका नन्हा-सा अंश ही कोशिकाओं में बदलता है।
\end{solution}

\begin{solution}{exo:b2:microbial-metabolism:9}
$L_p = \sqrt{2\times 1.2\times 10^{-9}\times 0.2/0.5} =
\sqrt{9.6\times 10^{-10}} = \qty{31}{\micro m}$। $c_0$ दुगुना करने पर
$L_p$ $\sqrt 2$ गुना हो जाता है: \qty{44}{\micro m}; $q$ चौगुना करने
पर वह आधा रह जाता है: \qty{15}{\micro m}।
\end{solution}

\begin{solution}{exo:b2:microbial-metabolism:10}
ऊपर से नीचे: शैवाल और सायनोजीवाणु सहित जल (ऑक्सीजनी
प्रकाशसंश्लेषण); लौह ऑक्सीकारियों की ज़ंग-रंगी पट्टी, जहाँ $\mathrm{Fe^{2+}}$
ऑक्सीजन से मिलता है; रंगहीन गंधक-ऑक्सीकारियों
(रसशिलापोषियों) का सफ़ेद परदा, जहाँ सल्फाइड ऑक्सीजन से मिलता है;
बैंगनी गंधक-जीवाणुओं की बैंगनी पट्टी (सल्फाइड पर \omterm{def:b2:microbial-metabolism:anoxygenic}{अनॉक्सीजनी प्रकाशसंश्लेषण});
हरे गंधक-जीवाणुओं की हरी पट्टी (वही, पर अधिक सल्फाइड और
कम प्रकाश सहते हुए); किण्वकारियों तथा सल्फेट अपचायकों की काली कीचड़।
(क) बैंगनी गंधक-जीवाणुओं को ऊपर से प्रकाश और नीचे से सल्फाइड, दोनों
चाहिए, इसलिए वे उसी अंतराफलक पर बैठते हैं जहाँ दोनों प्रवणताएँ कटती
हैं। (ख) सल्फेट अपचायक $\mathrm{H_2S}$ बनाते हैं, जो लोहे को काले FeS
के रूप में अवक्षेपित कर देता है। (ग) पदार्थ के लिए बंद: गंधक, कार्बन
और नाइट्रोजन बिना बाहर गए ऑक्सीकृत तथा अपचयित रूपों के बीच चक्कर काटते
हैं; ऊर्जा के लिए खुला: प्रकाश भीतर आता है, ऊष्मा बाहर जाती है, और
प्रकाश के बिना हर चक्र रुक जाता है।
\end{solution}

\begin{solution}{exo:b2:microbial-metabolism:11}
उत्पादन $pn$, हानि $kA$: $A^* = pn/k$। $A_c = 10^{-8}\times
6.02\times 10^{23} = 6.0\times 10^{15}$ अणु प्रति लीटर; $n_c =
kA_c/p = 5\times 10^{-4}\times 6.0\times 10^{15}/5 = 6\times 10^{11}$
प्रति लीटर, यानी $6\times 10^{8}$ प्रति मिलिलिटर। सघन संवर्धन
($10^{9}$) देहली से ऊपर है, समुद्री जल ($10^{5}$) उससे परिमाण की चार
कोटि नीचे। किसी \omterm{def:b2:microbial-metabolism:biofilm}{जैवफ़िल्म} में कोशिकाएँ आधात्री के प्रति मिलिलिटर
$10^{11}$ के घनत्व पर बैठती हैं, और आधात्री संकेत का निकलना धीमा कर
देती है (प्रभावतः $k$ घटाकर), इसलिए किसी नैनोलिटर में कुछ हज़ार
कोशिकाओं का सूक्ष्मनिवह कोरम तक पहुँच जाता है, जबकि वही कोशिकाएँ एक
लीटर में बिखरी हों तो कभी नहीं पहुँचतीं।
\end{solution}

\begin{solution}{exo:b2:microbial-metabolism:12}
नाइट्रोजन स्थिरीकरण (नाइट्रोजनेज़ केवल प्राक्केंद्रकियों के पास है):
रुक जाए तो जीवमंडल की यौगिक नाइट्रोजन \omterm{def:b2:microbial-metabolism:anaerobic}{विनाइट्रीकरण} से बह जाती
और उत्पादन दशकों में ढह जाता।
नाइट्रीकरण और \omterm{def:b2:microbial-metabolism:anaerobic}{विनाइट्रीकरण} (केवल \omterm{def:b2:unicellular-diversity:prokaryote}{प्राक्केंद्रकी}): रुक जाएँ तो
अमोनियम जमा होता और नाइट्रेट लुप्त, और नाइट्रोजन कभी वायु में नहीं
लौटती। \omterm{def:b2:microbial-metabolism:anaerobic}{मेथेनजनन} और \omterm{def:b2:microbial-metabolism:anaerobic}{अवायवीय श्वसन}
(केवल \omterm{def:b2:unicellular-diversity:prokaryote}{प्राक्केंद्रकी}): रुक जाएँ तो अनॉक्सी अवसादों, रूमेनों और दलदलों
का कार्बनिक पदार्थ बिना खनिजीकृत हुए जमा होता जाता और कार्बन चक्र से
बाहर निकल जाता; इसी तरह, ऑक्सीजनी प्रकाशसंश्लेषण एक जीवाणु-आविष्कार
था, और उसके बिना श्वसन के लिए ऑक्सीजन ही न होती। सुकेंद्रकी एक
\omterm{def:b2:unicellular-diversity:prokaryote}{प्राक्केंद्रकी} ग्रह में देर से जुड़ा और रासायनिक दृष्टि से सँकरा
जोड़ हैं।
\end{solution}

\begin{solution}{pb:b2:microbial-metabolism:1}
\textbf{1.} $-24\times 96.5\times (0.82 + 0.43) = \qty{-2900}{kJ/mol}$।
\textbf{2.} नाइट्रेट: $-24\times 96.5\times 1.17 = \qty{-2700}{kJ/mol}$;
सल्फेट: $-24\times 96.5\times 0.21 = \qty{-490}{kJ/mol}$।
\textbf{3.} ऑक्सीजन अधिक देती है, इसलिए जब तक ऑक्सीजन है तब तक
वायुजीवी कार्बनिक पदार्थ के लिए विनाइट्रीकारियों को हरा देते हैं, और
ऑक्सीजन नाइट्रेट रिडक्टेज़ों को दबा देती है; \omterm{def:b2:microbial-metabolism:anaerobic}{सल्फेट अपचयन} को अनॉक्सिता
\emph{और} नाइट्रेट की अनुपस्थिति दोनों चाहिए, इसलिए उसकी गंध का अर्थ है
कि टैंक कम वातित है और नाइट्रेट चुक चुकी है --- और $\mathrm{H_2S}$
विषैली है तथा कंक्रीट को खा जाती है।
\textbf{4.} $0.4\times 2900/50 = 23$ ATP; $0.4\times 2700/50 = 22$;
$0.4\times 490/50 = 4$।
\textbf{5.} लगभग छह गुना अधिक ($23/4$)।
\textbf{6.} प्रति ग्लूकोज़ दो ATP का अर्थ है लगभग \num{0.1} की उपज,
इसलिए क्रियाधार की ऊर्जा का \qty{90}{\%} उत्पादों (अम्ल,
ऐल्कोहॉल, $\mathrm{H_2}$) में ही रह जाता है, जिन्हें मेथेनजनक उतने ही
छोटे लाभ से मेथेन में बदल देते हैं; कम ऊर्जा, कम जैवभार, और कार्बन गैस
बनकर निकल जाता है।
\textbf{7.} $10^4\times 0.3 = \qty{3000}{kg}$ COD प्रतिदिन।
\textbf{8.} जैवभार $0.4\times 3000 = \qty{1200}{kg}$ COD; शेष
\qty{1800}{kg} COD ऑक्सीकृत होती है और उसे \qty{1800}{kg}
$\mathrm{O_2}$ चाहिए।
\textbf{9.} प्रतिदिन $1800/2 = \qty{900}{kWh}$।
\textbf{10.} $10^4\times 0.04 = \qty{400}{kg}$ नाइट्रोजन;
$4.57\times 400 = \qty{1830}{kg}$ $\mathrm{O_2}$।
\textbf{11.} स्थायी अवस्था में $\mu = 1/\theta \le \mu_{\max}$:
$\theta = 1/\mu_{\max} = \qty{2}{d}$ से नीचे बहाव।
\textbf{12.} $D = \qty{0.1}{d^{-1}}$: $S^* = 1\times 0.1/(0.5 - 0.1)
= \qty{0.25}{g/m^3}$।
\textbf{13.} $S^* = 0.1/(0.25 - 0.1) = \qty{0.67}{g/m^3}$, और बहाव
की सीमा चढ़कर \qty{4}{d} हो जाती है; संचालक को पंक की आयु बढ़ानी
पड़ेगी (कम पंक फेंके): $\theta = \qty{20}{d}$ पर $S^* =
0.05/0.2 = \qty{0.25}{g/m^3}$ फिर से।
\textbf{14.} प्रतिदिन $2.86\times 400 = \qty{1144}{kg}$ COD।
\textbf{15.} बहिःस्राव का अमोनियम $0.25\times 10^4 = \qty{2.5}{kg/d}$,
यानी भार का \qty{0.6}{\%}; शेष \qty{397}{kg} नाइट्रोजन
$\mathrm{N_2}$ बनकर निकल जाती है: \qty{99}{\%}।
\textbf{16.} बची: \qty{1144}{kg} $\mathrm{O_2}$। कुल माँग:
प्रतिदिन $(1800 - 1144) + 1830 = \qty{2490}{kg}$ $\mathrm{O_2}$।
\textbf{17.} प्रतिदिन $1200\times 0.6\times 0.35 = \qty{252}{m^3}$
मेथेन।
\textbf{18.} $252\times 36 = \qty{9070}{MJ} = \qty{2520}{kWh}$;
बिजली $0.35\times 2520 = \qty{880}{kWh}$ प्रतिदिन।
\textbf{19.} वातन प्रतिदिन $2490/2 = \qty{1245}{kWh}$: मेथेन उसका
$880/1245 = \qty{71}{\%}$ ढँक लेती है।
\textbf{20.} $D_e\,c'' = q$, इसलिए $c = c_0 + ax + qx^2/2D_e$;
$c(L_p) = 0$ और $c'(L_p) = 0$ के साथ: $a = -qL_p/D_e$ और $c =
(q/2D_e)(L_p - x)^2$, जिसमें $L_p^2 = 2D_ec_0/q$।
\textbf{21.} $L_p = \sqrt{2\times 1.5\times 10^{-9}\times 0.2/0.3} =
\sqrt{2\times 10^{-9}} = \qty{45}{\micro m}$।
\textbf{22.} $45/300 = \qty{15}{\%}$ वायवीय। उसके नीचे विनाइट्रीकारी
नाइट्रीकारी सतही परत से नीचे विसरित होती नाइट्रेट को, भीतर विसरित होते
कार्बनिक पदार्थ से, $\mathrm{N_2}$ तक अपचयित करते हैं।
\textbf{23.} एक ही फ़िल्म में ऑक्सीजन प्रवणता एक वायवीय परत
(नाइट्रीकरण) को एक अनॉक्सी परत (\omterm{def:b2:microbial-metabolism:anaerobic}{विनाइट्रीकरण}) के ऊपर खड़ा कर
देती है, और दोनों कुछ दहाई माइक्रोमीटर ही दूर होती हैं; वातित टैंक पूरा
ऑक्सीजनी है, इसलिए \omterm{def:b2:microbial-metabolism:anaerobic}{विनाइट्रीकरण} के लिए अलग अनॉक्सी टैंक चाहिए।
\textbf{24.} क्लोरीन गहरी कोशिकाओं तक पहुँचने से पहले ही आधात्री के
बहुलकों से अभिक्रिया करके ख़र्च हो जाती है; गहरी कोशिकाएँ भूखी,
धीमी-बढ़त वाली या प्रसुप्त हैं और कहीं कम संवेदी; केवल सतही परत मरती है
और फ़िल्म नीचे से फिर उग आती है।
\textbf{25.} ऑक्सीजन माँग प्रतिदिन \qty{2490}{kg}; मेथेन प्रतिदिन
\qty{252}{m^3}; मेथेन वातन बिजली का \qty{71}{\%} जुटाती है
(\qty{1245}{kWh} में से \qty{880}{kWh})।
\end{solution}
